\chapter{PELAKSANAAN PENELITIAN}\label{pelaksanaan-penelitian}
% Structure follows template-kampus/Bab3.tex. The guideline also calls this chapter
% "Metodologi Penelitian" (alir penelitian, algoritma, setup) - confirm the title with the advisor.

\section{Alat dan Bahan Penelitian}\label{alat-bahan}
% Spesifikasi perangkat keras/lunak dan bahan (dataset) yang digunakan.
[ISI NANTI]

\section{Tata Laksana Penelitian}\label{tata-laksana}
% Rangkaian logis penyelesaian masalah: langkah kerja dan/atau algoritma penelitian.

\subsection{Alur Penelitian}
[ISI NANTI]

\subsection{Dataset}
[ISI NANTI]

\subsection{Praproses Data}
[ISI NANTI]

\subsection{Ekstraksi Fitur}
[ISI NANTI]

\subsection{Model Klasifikasi}
[ISI NANTI]

\subsection{Skema Pelatihan dan Pembagian Data}
[ISI NANTI]

\section{Rencana Analisis Hasil}\label{rencana-analisis}
% Bagaimana hasil diolah sesuai tujuan: ringkasan tabel/gambar/statistik, uji perbedaan/korelasi,
% evaluasi model dan Explainable AI.
[ISI NANTI]
